% =============================================================================
% thesis_fa_digits.tex؛ ارقام در متن RTL فارسی با \pnum{...}
% =============================================================================
% \pnum{0.746} → عدد فارسیِ امن داخل متن فارسی.
% عدد در منبع ASCII/LTR می‌ماند، در خروجی به رقم فارسی نگاشت می‌شود، و داخل جزیرۀ
% LTR قرار می‌گیرد؛ این مسیر از وارونگی نمونه‌هایی مثل ۰۰۰،۲ جلوگیری می‌کند.
% انگلیسی LTR: \lr{AUROC 0.746} یا \cnum{0.746} (= همان \lr{...})
% ریاضی، verbatim، listings دست‌نخورده بمانند.
% =============================================================================

\ExplSyntaxOn
\cs_new:Npn \thesis_pnum_map:n #1
  {
    \tl_map_inline:nn {#1}
      {
        \str_case:nnF { ##1 }
          {
            {0}{۰}
            {1}{۱}
            {2}{۲}
            {3}{۳}
            {4}{۴}
            {5}{۵}
            {6}{۶}
            {7}{۷}
            {8}{۸}
            {9}{۹}
            {.}{\ThesisDecimalSep}
            {,}{\ThesisThousandsSep}
          }
          { ##1 }
      }
  }
\NewDocumentCommand \pnum { m }
  {
    \mbox{\beginL{\thesis_pnum_map:n {#1}}\endL}
  }
\NewDocumentCommand \prange { m m }
  {
    \mbox{\beginL{\thesis_pnum_map:n {#1}{\ThesisLatinPunct -}\thesis_pnum_map:n {#2}}\endL}
  }
\ExplSyntaxOff
